%% paper/sections/07_results.tex
%% §7 Results — 분류군별 대표 검증 framing + 새 Table II 포함

\section{측정 결과}
\label{sec:results}

본 절은 \S\ref{subsec:hypotheses} 의 세 가설 (H1: 교차 도메인 중첩, H2: 박자 정렬, H3: 분류군별 실패 메커니즘) 의 경험적 검증 결과를 보고한다.
\algname 의 5단계 hyperparameter 의 최종값 (Table~\ref{tbl:metronome-hyperparam}) 은 이 측정들에서 직접 도출된다.

%% tables/tbl_vs_baseline.tex
%% Table I — 5 분류군 대표 기법의 baseline 대비 효과 (table* + resizebox)

\begin{table*}[t]
\centering
\caption{5 분류군 대표 기법의 baseline 대비 효과 ($\Delta$ tps / wall, 1-run 측정). \algname 의 단일 통합 구성 (NUMA pin $\oplus$ branchy aux) 이 유일하게 유의 임계점 $+5\,\%$ 에 도달.}
\label{tbl:vs-baseline}
\scriptsize
\setlength{\tabcolsep}{4pt}
\resizebox{\textwidth}{!}{%
\begin{tabular}{@{}clp{4.6cm}rrll@{}}
\toprule
\textbf{Rank} & \textbf{Class} & \textbf{Representative} & \textbf{$\Delta$ tps} & \textbf{$\Delta$ wall} & \textbf{Domain} & \textbf{Outcome} \\
\midrule
\textbf{1} & E$\oplus$B & \algname\ (NUMA $\oplus$ branchy 100ms) & $\mathbf{+5.28\%}$ & $\mathbf{-5.40\%}$ & cross-domain        & \textbf{threshold} \\
2          & E          & branchy aux $\times$ 100\,ms             & $+5.28\%$           & $-5.40\%$           & process work        & cell $+12.04\%$ \\
3          & B$\oplus$E & NUMA $\oplus$ regular softmax aux        & $+2.83\%$           & $-2.52\%$           & cross-domain        & near-linear \\
4          & B          & NUMA pin (dual instance)                 & $+1.84\%$           & $-1.71\%$           & system              & robust (m-run $+2.24\%$) \\
5          & E          & regular aux $\times$ 100\,ms             & $+1.84\%$           & $-1.71\%$           & process work        & noise \\
6          & E          & async tokenizer $\times$ 20\,ms          & $+1.66\%$           & $-1.77\%$           & process work        & m-run 부호 반전 \\
7          & A          & AMX matmul drop-in                       & $+0.26\%$           & $+0.00\%$           & critical-path       & H2D/D2H sync \\
8          & B          & cgroup + hugepages + taskset             & $-0.39\%$           & $+0.41\%$           & system              & same-domain stack \\
9          & C          & task-pool dummy fill                     & $-1.75\%$           & $+1.92\%$           & sub-step            & lock contention \\
10         & A          & AMX prefill assist                       & $-2.08\%$           & $+2.12\%$           & critical-path       & drop-in 불가 \\
11         & D          & AMX draft proxy real integration         & $-2.77\%$           & $+2.85\%$           & NEW workload        & invasive 부재 \\
12         & A          & 4-method Jacobi router                   & $-94.20\%$          & ---                 & critical-path       & queue collapse \\
\bottomrule
\end{tabular}%
}
\smallskip
\par\noindent\scriptsize 측정: Qwen 2.5-32B + TP=4$\times$2 + H100$\times$8, 500 prompts $\times$ \texttt{max\_tokens}=32, 3-mix workload, 1-run. Class: A=직접 커널, B=시스템 격리, C=sub-step, D=보조 워크로드, E=동시 보조. m-run = multi-run.
\end{table*}


\subsection{핵심 결과 1: \algname 의 임계점 도달}

Table~\ref{tbl:vs-baseline} 의 1행은 \algname 의 단일 통합 구성 (NUMA pin $\oplus$ branchy auxiliary 100\,ms) 의 측정 결과이다.
3-mix 평균 $\Delta = +5.28\,\%$, wall-clock 단축 $-5.40\,\%$ 로 유의 임계점에 도달한다.
단일 cell (balanced workload) 에서 최대 $\Delta = +12.04\,\%$ 까지 확장된다.

비교 대상은 \algname 의 두 구성 요소인 NUMA pin 단독 (4행, $+1.84\,\%$) 과 branchy auxiliary 단독 (2행, $+5.28\,\%$)이다.
\algname 의 통합 구성이 두 단일 기법보다 정합적 (consistent) 으로 우월하며, 이는 \textsc{Chord} 단계의 교차 도메인 화음이 단순 합을 넘어 \emph{시스템적 효과}를 만들어냄을 시사한다.

\subsection{핵심 결과 2: 분류군 대표 검증 (H3)}

%% tables/tbl_class_representatives.tex
%% Table II — 분류군별 대표 기법 + 메커니즘 (table* + resizebox, multirow 제거)

\begin{table*}[t]
\centering
\caption{5 분류군 대표 기법의 단독 효과와 실패/성공 메커니즘. 각 분류군은 자원의 도메인에 따라 정의되며, 동일 분류군 내 기법들은 공통 메커니즘으로 실패 또는 성공한다.}
\label{tbl:class-representatives}
\scriptsize
\setlength{\tabcolsep}{4pt}
\resizebox{\textwidth}{!}{%
\begin{tabular}{@{}cp{4.3cm}rlp{6.0cm}@{}}
\toprule
\textbf{Class} & \textbf{Representative} & \textbf{$\Delta$ (1-run)} & \textbf{Verdict} & \textbf{Mechanism} \\
\midrule
A & AMX matmul drop-in              & $+0.26\%$  & noise         & H2D/D2H sync overhead 누적이 microbench gain 상쇄 \\
A & 4-method Jacobi router          & $-94.20\%$ & catastrophic  & LM-head 1.81\,s $\times$ 165 chat req 직렬 큐 collapse \\
B & NUMA pin (dual instance)        & $\mathbf{+1.54\%}$ ($\uparrow$ multi $+2.24\%$) & \textbf{robust} & local/remote 2.1$\times$ latency 비대칭의 hardware-level 회수 \\
B & 4-OS-lever stack                & $-0.03\%$  & destructive   & 동일 도메인 stacking $\rightarrow$ context-switch + cache thrash \\
C & task-pool fill (actual fire)    & $-1.75\%$  & catastrophic  & step-internal lock contention 으로 critical path stall \\
D & AMX draft proxy (Qwen 0.5B)     & $-2.77\%$  & failure       & microbench 0.524\,ms, invasive integration 부재로 e2e 변환 실패 \\
E & branchy $\times$ 100\,ms (cellB) & $\mathbf{+5.28\%}$ (cell $+12.04\%$) & \textbf{threshold} & step-boundary 정렬 + prefetcher 억제의 결합 효과 \\
E & regular $\times$ 10\,ms (cellA)  & $+0.98\%$  & noise         & 비정수배 alignment 로 critical-path 침범 \\
\bottomrule
\end{tabular}%
}
\smallskip
\par\noindent\scriptsize 본 분류군 대표 검증은 가설 H1--H3 (\S\ref{subsec:hypotheses}) 의 정당화 근거. 분류군 B 와 E 의 cross-domain pair (NUMA $\oplus$ branchy) 가 \algname 의 \textsc{Chord} + \textsc{Accent} 단계 design 의 직접 backing (Table~\ref{tbl:superposition}).
\end{table*}


Table~\ref{tbl:class-representatives} 는 5 분류군의 대표 기법과 각 분류군의 실패/성공 메커니즘이다.

\textbf{분류군 A (직접 커널 대체).}
AMX matmul drop-in 은 microbench 1.5--3$\times$ 의 speedup 에도 e2e $\Delta = +0.26\,\%$ noise 에 머문다.
극단적 사례로 Jacobi LM-head ($\sim$1.81\,s) 의 직접 대체 시도는 165 chat 요청의 직렬 큐 collapse 로 $-94.20\,\%$ catastrophic.
\textbf{메커니즘}: critical-path 위 H2D/D2H 동기화 비용 (PCIe transfer + CUDA stream sync $\sim$5.5\,ms) 이 microbench gain 을 상쇄.

\textbf{분류군 B (시스템 격리).}
NUMA pin 단독은 1-run $+1.84\,\%$ 에서 multi-run mean $+2.24\,\%$ 으로 일관된 양성 신호를 유지한다 (Table~\ref{tbl:multirun}).
반면 cgroup + hugepages + taskset 동시 stacking 은 $-0.39\,\%$ destructive.
\textbf{메커니즘}: NUMA distance 비대칭 (local 10 vs.\ remote 21) 은 hardware-level 일관 효과를 제공.
같은 도메인 stacking 은 OS context-switch / cache contention 누적.

\textbf{분류군 C (sub-step 스케줄링).}
task-pool 기반 dummy fill 의 실제 fire 시 $-1.75\,\%$, spec-decode backend 에서 $-14$\,\% 까지 catastrophic.
\textbf{메커니즘}: step 내부의 sub-phase 경계 (35--44\,ms 영역) 에서 task-pool lock contention 이 critical path stall 을 유발.

\textbf{분류군 D (보조 워크로드).}
AMX draft head 의 isolated microbench 는 $0.524$\,ms 의 매우 빠른 결과를 보이지만, e2e proxy 측정에서는 $-2.77\,\%$.
\textbf{메커니즘}: spec decode integration 의 invasive 변경 (rejection sampling tree path, KV cache layout) 부재로 microbench-to-e2e 변환이 자동으로 일어나지 않음.

\textbf{분류군 E (동시 보조 작업).}
$\tcpu = 100$\,ms branchy pattern 이 $+5.28\,\%$ 로 유의 임계점 도달.
대조적으로 $\tcpu = 10$\,ms regular pattern 은 $+0.98\,\%$ noise.
\textbf{메커니즘}: 박자 정렬 + 분기 패턴의 prefetcher 억제 (\S\ref{sec:mechanism}).

\subsection{핵심 결과 3: 가설 H1 (교차 도메인 중첩) 검증}

%% tables/tbl_superposition.tex
%% Table III — Superposition (table* + resizebox)

\begin{table*}[t]
\centering
\caption{Theorem~\ref{thm:superposition} 의 경험적 검증. 도메인 간섭 비용 $\varepsilon$ 가 동일 도메인 stacking 에서 1.18\,pp, 교차 도메인에서 0.55\,pp 로 약 2배 격차를 보이며, 3+ 동일 도메인 stack 은 superlinear destructive (4.56\,pp).}
\label{tbl:superposition}
\scriptsize
\setlength{\tabcolsep}{4pt}
\resizebox{\textwidth}{!}{%
\begin{tabular}{@{}lccccc@{}}
\toprule
\textbf{Stack composition} & \textbf{Domain relation} & \textbf{Linear sum} & \textbf{Actual} & \textbf{Residual $\varepsilon$} & \textbf{Verdict} \\
\midrule
NUMA pin $\oplus$ softmax aux        & \textbf{cross} (B$\oplus$E) & $+3.38\%$ & $+2.83\%$ & $\mathbf{-0.55\,pp}$ & near-linear \\
cgroup $\oplus$ taskset              & same (B$\oplus$B)           & $+0.06\%$ & $-0.39\%$ & $-0.45\,pp$          & destructive \\
4-OS-lever (cgroup+hp+taskset+NUMA)  & same (B$^{4}$)               & $+1.15\%$ & $-0.03\%$ & $\mathbf{-1.18\,pp}$ & destructive \\
Triple-OS-lever                      & same (B$^{3}$)               & $+1.10\%$ & $-2.83\%$ & $\mathbf{-3.93\,pp}$ & superlinear destructive \\
\bottomrule
\end{tabular}%
}
\smallskip
\par\noindent\scriptsize Empirically derived: $\varepsilon_{\text{cross}} \!\approx\! 0.55\,$pp $\ll \varepsilon_{\text{same}}^{(2)} \!\approx\! 1.18\,$pp $\ll \varepsilon_{\text{same}}^{(3+)} \!\approx\! 4.56\,$pp.
\end{table*}


Table~\ref{tbl:superposition} 는 Theorem~\ref{thm:superposition} 의 직접 검증 결과이다.

\textbf{Cross-domain pair (NUMA $\oplus$ regular softmax aux)}: 단일 효과의 단순 합 $+3.38\,\%$ 에 대해 실측 $+2.83\,\%$, residual $-0.55$\,pp.
이는 near-linear superposition 의 강한 신호이며, \algname 의 \textsc{Chord} 단계 design 의 직접 근거이다.

\textbf{Same-domain stack (4-OS-lever)}: 단순 합 $+1.15\,\%$ 예측, 실측 $-0.03\,\%$, residual $-1.18$\,pp.
3-기법 stacking 으로 확장 시 residual 이 $-3.93$\,pp 까지 발산하여 superlinear destructive 의 강한 신호를 보인다.

이 결과는 Theorem~\ref{thm:superposition} 의 경험적 정당화를 제공한다.

\subsection{핵심 결과 4: 가설 H2 (박자 정렬) 검증}

분류군 E 의 2$\times$2 ablation 결과 (Fig.~\ref{fig:pattern-cycle-grid}) 는 다음 관찰을 제공한다.
$\tcpu = 100$\,ms ($k = 2.5$) $\times$ branchy 패턴이 $+5.28\,\%$ 로 임계점 도달.
$\tcpu = 10$\,ms ($k = 0.25$) regular 는 $+0.98\,\%$, $\tcpu = 20$\,ms ($k = 0.5$) regular 는 1-run $+1.66\,\%$ 이나 multi-run 부호 반전.

\emph{$k$ 가 정수배 (2 또는 3) 에 도달할 때만} 정렬 이득이 나타나며, $k < 1$ 또는 비정수배는 비정상적 정렬로 critical path 침범 가능성이 높다.
이는 Theorem~\ref{thm:alignment} 의 경험적 검증이다.

\subsection{Multi-run 분산의 함의}

%% tables/tbl_multirun.tex
%% Table IV — Multi-run variance: cold-start outlier 영향과 warm-only signal

\begin{table}[t]
\centering
\caption{핵심 후보 기법의 multi-run 분산 분석. cold-start outlier (1회) 분리 시 warm-only 신호가 1-run 측정과 다르게 부호 또는 magnitude가 변화하는 경우가 발생함.}
\label{tbl:multirun}
\small
\begin{tabular}{@{}lrrrl@{}}
\toprule
\textbf{Technique} & \textbf{1-run} & \textbf{3-run mean} & \textbf{warm-only} & \textbf{Verdict} \\
\midrule
NUMA pin           & $+1.54\%$ & $\mathbf{+2.24\%}$ & $\mathbf{+3.13\%}$ & \textbf{robust positive} \\
softmax aux 100ms  & $+1.84\%$ & $+0.53\%$          & $+0.81\%$          & noise floor \\
tokenize aux 20ms  & $+1.66\%$ & $\mathbf{-5.96\%}$ & $\mathbf{-6.40\%}$ & \textbf{retract} \\
\bottomrule
\multicolumn{5}{l}{\footnotesize 1-run 측정에서 부호가 양이라도 multi-run 평균이 부호 반전을 보이는} \\
\multicolumn{5}{l}{\footnotesize 사례(tokenize 20\,ms)가 존재하므로, magnitude $<3\%$ 의 단일 측정은} \\
\multicolumn{5}{l}{\footnotesize 신뢰도가 낮다. NUMA pin은 hardware 비대칭에 기반하여 robust.} \\
\end{tabular}
\end{table}


Table~\ref{tbl:multirun} 의 핵심 관찰은 NUMA pin 만 multi-run 평균에서 robust 양성 (1-run $+1.54\,\%$ $\rightarrow$ steady-state $+3.13\,\%$) 을 유지한다는 점이다.
async tokenizer 20\,ms 는 1-run $+1.66\,\%$ 의 양성 신호가 multi-run 에서 $-5.96\,\%$ 로 부호 반전한다.

이는 LLM 추론 시스템의 multi-run 측정 best practice 를 시사한다.
1-run 측정에서 magnitude $< 3\,\%$ 인 신호는 noise 와 분리되지 않으므로, hardware-level 이론 backing 이 없는 단일 신호는 단독으로 채택하지 않는다.

\subsection{\algname Hyperparameter 의 도출}

%% tables/tbl_metronome_hyperparam.tex
%% Table V — METRONOME 단계별 hyperparameter 와 그 경험적 도출 근거

\begin{table}[t]
\centering
\caption{\algname 의 단계별 hyperparameter 와 경험적 도출 근거.}
\label{tbl:metronome-hyperparam}
\small
\begin{tabular}{@{}lll@{}}
\toprule
\textbf{Stage} & \textbf{Hyperparameter} & \textbf{값 / 근거} \\
\midrule
\textsc{Tempo}    & probe sample size $n$       & $n = 10$ steps \\
\textsc{Tempo}    & $\tgpu$ (측정값)            & $\approx 40$\,ms (Qwen 2.5-32B + TP=4) \\
\midrule
\textsc{Chord}    & $L_{\text{primary}}$        & NUMA pin (domain = system) \\
\textsc{Chord}    & $L_{\text{secondary}}$      & auxiliary work (domain = work) \\
\textsc{Chord}    & domain constraint           & $\dom(L_{\text{primary}}) \neq \dom(L_{\text{secondary}})$ \\
\midrule
\textsc{Meter}    & integer multiplier $k$      & $k \in \{2, 3\}$, default 2.5 \\
\textsc{Meter}    & $\tcpu^{\min}$              & 100\,ms \\
\textsc{Meter}    & $\tcpu$ (도출값)            & $\max(2\tgpu, 100\,\mathrm{ms}) = 100$\,ms \\
\midrule
\textsc{Accent}   & PCIe pressure threshold $\tau$ & 0.5 (정규화) \\
\textsc{Accent}   & 본 환경 pattern               & BRANCHY (8-GPU TP $\rightarrow$ PCIe 압박) \\
\midrule
\textsc{Ritardando} & drift threshold $\delta_{\max}$ & 0.20 \\
\textsc{Ritardando} & monitor interval            & per-step \\
\bottomrule
\multicolumn{3}{l}{\footnotesize 모든 hyperparameter는 \S\ref{sec:results}의 24-기법 ablation으로부터 경험적 도출.} \\
\end{tabular}
\end{table}


Table~\ref{tbl:metronome-hyperparam} 는 위 검증 결과로부터 도출된 \algname 의 단계별 hyperparameter 이다.
\textsc{Tempo} 의 $n=10$ 은 $\tgpu$ 측정의 분산 안정화 sample size, \textsc{Meter} 의 $\tcpu = 100$\,ms 는 가설 H2 의 ablation 최적, \textsc{Accent} 의 BRANCHY 는 8$\times$ TP 환경의 PCIe pressure 측정값 기반이다.

본 hyperparameter 는 본 측정 환경에 특화되지만, \textsc{Tempo} 자동 측정과 \textsc{Ritardando} drift 재조정이 다른 model/workload 환경으로의 일반화를 부분적으로 지원한다.
